\section{Methodology overview (before the equations)}
\label{sec:method-overview}

The Methodology section of the talk has two layers: (i)~what data and system pieces you use,
(ii)~the mathematics of each pipeline stage. This chapter covers layer~(i). The next chapter
goes to the mathematical core.

\subsection{Chief complaints as input (slide 11)}

\begin{definition}[Chief complaint \(X\)]
The primary reason the patient presents, in their own words, as unstructured free text at intake.
It is \emph{not} a diagnosis, not TEWS, and not a full medical history.
\end{definition}

\begin{whybox}
Forces clarity on the input variable. Everything later is a function of \(X\).
\end{whybox}

\paragraph{Why English only?}
BioBERT's WordPiece vocabulary and the Triagegeist training pairs are English.
The prototype does not claim Twi or mixed-language support.

\paragraph{Running example.}
At \examplehospital{} we reuse: \runningcomplaint{}
Walk this sentence through gate \(\to\) tokens \(\to\) colour when an examiner asks for an example.

\begin{saybox}
``Our only input is the English chief complaint text \(X\). Vitals are the nurse's job in parallel.''
\end{saybox}

\subsection{Training data provenance (slide 12)}

\begin{whybox}
Stops the ``where did the numbers come from?'' question. Two layers must never be confused.
\end{whybox}

\paragraph{Layer A --- Triagegeist (metrics source).}
Synthetic ED benchmark \citep{triagegeist2025}.
Merge \texttt{train.csv} + \texttt{chief\_complaints.csv} \(\to\) \(N=80{,}000\) English pairs
(complaint text, acuity label).
Stage~3 project fine-tuning and the 90/10 holdout use these pairs.
\textbf{All reported accuracy / F1 / recall numbers come from Layer~A.}

\paragraph{Layer B --- KNUST visit examples.}
Written after a site visit to \examplehospital{} (workflow, SATS colours, departments).
Used for pathway demos and qualitative checks.
\textbf{Not} merged into the 80k training set---so they do not inflate metrics.

\begin{saybox}
``Eighty thousand Triagegeist pairs train and evaluate the classifier. KNUST examples shape pathways and demos; they are not in the training metrics.''
\end{saybox}

\subsection{NLP pipeline overview (slide 13)}

\begin{whybox}
One diagram for the whole story before equations. Examiners should see the chain once.
\end{whybox}

Pipeline stages in order:
\[
  X \;\to\; \text{clinical gate} \;\to\; \text{tokenise} \;\to\; \text{embed (768-d)}
  \;\to\; \text{BioBERT (\(L=12\))} \;\to\; \text{softmax (5 classes)}
  \;\to\; \text{SATS colour} \;\to\; P(C).
\]

\begin{itemize}
  \item Gate rejects non-clinical text (greetings) before BioBERT runs.
  \item Softmax produces five probabilities that sum to \(1\).
  \item Chosen level \(\hat{c}\): the class with the \emph{highest} probability.
  \item Fixed map \(C_{\mathrm{NLP}}=f_{\mathrm{SATS}}(\hat{c})\), then pathway card \(P(C_{\mathrm{NLP}})\).
\end{itemize}

\begin{saybox}
``Gate first, then BioBERT, then pick the most likely acuity level, map it to a colour, and attach a pathway card.''
\end{saybox}

\subsection{Pathways at \examplehospital{} (slide 14)}

\begin{whybox}
Connects colours to destinations and timers---the clinical product, not only a class label.
\end{whybox}

\begin{center}
\small
\begin{tabular}{@{}lll@{}}
\toprule
Colour & Destination (summary) & Timer \\
\midrule
Gate reject & Re-prompt for clinical text & n/a \\
Red & Resuscitation / immediate & Immediate \\
Orange & Acute / high-dependency & \(\approx 10\) min \\
Yellow & Urgent waiting / OPD urgent & \(\approx 60\) min \\
Green & Routine OPD / digital queue & Longer \\
\bottomrule
\end{tabular}
\end{center}

Seventeen hospital units were mapped (clinical + support). Pathway cards carry colour, destination, timer, and escalation.

\begin{saybox}
``A colour alone is not enough. The card says where to go and how soon.''
\end{saybox}

\subsection{System architecture (slide 15)}

\begin{whybox}
Shows the software path and that nurses remain in the loop.
\end{whybox}

Solid path: patient/attendant \(\to\) web UI \(\to\) \texttt{/predict} API \(\to\) gate \(\to\) BioBERT
\(\to\) pathway card \(\to\) department routing.

Dashed path: nurse TEWS in parallel into routing. Staff can override.

\begin{saybox}
``The chatbot is an assistive API in a web UI. TEWS and clinical override stay with the nurse.''
\end{saybox}
